\section{问题分析}

本题以镇龙乡标准测试场景为背景，围绕“运输组批—多架次调度—通信中继协同—任务分区”四个递进环节展开。四个子问题在物理规则（附录 2 时间与能耗）、通信规则（附录 3 链路判定）、对象编号与计算口径上完全统一，并逐层继承数据与中间结果。总体技术路线如图~\ref{fig:roadmap} 所示。

\begin{figure}[htbp]
  \centering
  \includegraphics[width=0.86\textwidth]{roadmap.png}
  \caption{全文技术路线图}
  \label{fig:roadmap}
\end{figure}

\subsection{问题一的分析}

问题一在不考虑实体无人机与共享电池调度的前提下，研究单服务区直接往返（$O01\to S_i\to O01$）的运输能力与货箱组批。其核心是两个耦合的子任务：一是确定各机型在各服务区的\textbf{最大安全载荷}，二是在载荷、体积与返航安全能量余量约束下完成\textbf{货箱组批}并优化。

最大安全载荷问题本质是由返航安全能量余量约束隐式定义的临界载荷求解：随载荷增大，等效航程缩短、单架次能耗上升，当能耗恰好触及安全余量边界时得到能量受限载荷，与机型最大载货质量取较小值即为最大安全载荷。由于能耗对载荷单调，可用一维求根（brentq）稳定求解。组批问题是典型的\textbf{装箱/集合划分}问题：货箱不可拆分、每箱仅安排一次，需同时满足质量与体积两类背包约束，并以架次数、总能耗、累计作业时间为多目标。该问题规模可控（单服务区货箱数有限），适合先枚举生成候选批次池，再用多目标 MILP 精确求解。本文选用\textbf{线性加权法}与\textbf{理想点法（Chebyshev 妥协规划）}两种标量化方法交叉验证。

\subsection{问题二的分析}

问题二放开“单点往返”限制，每个架次可多点访问，并引入实体机队（A/B/C 三型 8 架）、共享电池数量、两阶段充电周转以及医疗物资期望送达、首批保障截止等时限约束，是一个\textbf{带资源约束与时间窗的多行程车辆路径问题（MTVRP）}。决策空间同时包含组批、访问顺序、机型指派、实体机与电池指派、架次开始时刻，维度高且强耦合，精确求解难以在合理时间内完成。因此本文采用\textbf{分组编码 NSGA-II} 进行全局多目标优化，以“加权迟到量—完工时间—总能耗—架次数”为目标，并用贪心构造解作为基线对照，刻画“零违约但架次多”与“架次少但易违约”的权衡关系。

\subsection{问题三的分析}

问题三在问题二基础上叠加\textbf{连续通信约束}：运输无人机在爬升、巡航、下降及投送全程必须保持与固定网关 $G01$ 的直连，或由一架中继无人机提供“运输—中继—$G01$”两段链路。通信状态判定依赖附录 3 的自由空间损耗、双向链路预算与 $30\,$m DEM 视线遮挡。求解需先建立\textbf{通信服务状态模型}，将航段按时间/空间采样判定直连是否可用；对存在缺口的航段，将中继悬停位置离散化为“缺口腿分数点 $\times$ 悬停高度档”候选网格，再\textbf{联合优化}运输方案与中继悬停位置、飞行海拔、服务时段及能源组件分配。目标扩展为运输与中继两类无人机的及时性、联合完工时间、总能耗与总架次。

\subsection{问题四的分析}

问题四以问题三的联合调度方案为不变基础，研究\textbf{任务分区}对资源配置的影响。约束是：保持组批、访问顺序、运输—中继任务与通信关系不变，多站架次所涉服务区须整组同划，各组独立执行、资源不得跨组调配。因此可先将“必须同组”的服务区经 \textbf{must-link 收缩}为原子任务块，再在原子块上枚举 $K=2,3$ 的所有可行分区。由于原子块数量不大，$K=2$ 可\textbf{完全枚举}精确求解，$K=3$ 辅以\textbf{局部搜索}。评价维度为资源配置规模、冗余、组间工作量均衡系数以及与现有库存的缺口。
